% Tabla C-T1. Generada automáticamente; no editar a mano.
% Requiere \usepackage{booktabs} (y graphicx si la tabla es ancha).
\begin{table}[htbp]
    \centering
    \caption{Pares de planetas adyacentes: estadística descriptiva de la separación orbital y de la desigualdad de tamaños, global y por tipo de par.}
    \label{tab:C-T1}
    \small
    \resizebox{\textwidth}{!}{%
    \begin{tabular}{lrrrrrrr}
    \toprule
    Variable & $N$ & Mediana & $Q_1$ & $Q_3$ & IQR & Mín. & Máx. \\
    \midrule
    period\_ratio & 898 & 2.166 & 1.720 & 3.121 & 1.400 & 1.038 & 54.083 \\
    $\ln(P_{\rm ext}/P_{\rm int})$ & 898 & 0.773 & 0.543 & 1.138 & 0.596 & 0.038 & 3.990 \\
    radius\_ratio & 898 & 1.150 & 0.949 & 1.480 & 0.531 & 0.412 & 3.675 \\
    $U_R$ & 898 & 0.222 & 0.095 & 0.438 & 0.343 & 0.000 & 1.302 \\
    abs\_log\_radius\_ratio en par\_rocoso & 132 & 0.082 & 0.038 & 0.157 & 0.119 & 0.000 & 0.464 \\
    abs\_log\_radius\_ratio en par\_mixto & 270 & 0.531 & 0.372 & 0.681 & 0.309 & 0.037 & 1.302 \\
    abs\_log\_radius\_ratio en par\_volatil & 496 & 0.165 & 0.080 & 0.303 & 0.223 & 0.000 & 0.876 \\
    \bottomrule
    \end{tabular}
    }
    \par\smallskip
    {\footnotesize\textit{Nota:} 898 pares en 615 sistemas. Fracción de pares con el planeta exterior mayor: 0.688.}
\end{table}
